\begin{multicols}{2}
\section{Prompts for Fuzzy PRD Semantic Equivalence Judge}
\label{appendix:fuzzy-prd-semantic}
\end{multicols}

\begin{promptbox}[Fuzzy PRD: Semantic Equivalence Judge]
\textbf{\#\# Role}

You are an expert benchmark-data auditor. Your job is to judge whether a generated fuzzy task package preserves the original detailed PRD well enough to support a paper claim about semantic equivalence, while separately identifying execution-contract risks that could affect reuse of the same tests.

\textbf{\#\# Inputs}
\begin{itemize}
\item Detailed PRD: \texttt{start.md}
\item Generated task package: JSON containing \texttt{fuzzy\_prd} and \texttt{oracle\_data}
\end{itemize}

Evaluate \texttt{fuzzy\_prd + oracle\_data} jointly rather than \texttt{fuzzy\_prd} alone.

\textbf{\#\# Key outputs}
\begin{itemize}
\item \texttt{task\_semantic\_similarity} in $[0,1]$
\item \texttt{task\_semantic\_equivalent}
\item \texttt{execution\_contract\_risk} in \{low, medium, high\}
\item \texttt{likely\_to\_fail\_same\_tests}
\item \texttt{paper\_support\_label}
\end{itemize}

Set \texttt{task\_semantic\_equivalent=true} when similarity is at least 0.90 and there is no material task-level contradiction that would change the implementation or behavioral tests.
\end{promptbox}






\begin{multicols}{2}
\section{Prompts for Fuzzy L3 PRD Construction}
\label{appendix:fuzzy-prd-easy-prompts}
\end{multicols}

\begin{promptbox}[Fuzzy L3 Generation: Constraint Extraction]
\textbf{\#\# Role}

You are a senior engineer performing precision analysis of a PRD.

\textbf{\#\# Task}

Read the detailed PRD together with all supporting files and produce a JSON object with key \texttt{"constraints"} containing every distinct behavioral constraint. The extraction must cover:
\begin{itemize}
\item mapping and branching behavior,
\item exact output field names and sentinel values,
\item action/dispatch contracts,
\item tested edge cases,
\item optional parameters and defaults,
\item validation and error-rejection rules.
\end{itemize}

Each item in \texttt{constraints} has the form:

\texttt{\{}\\
\texttt{~~"id": "C001",}\\
\texttt{~~"feature": "<short feature name>",}\\
\texttt{~~"input\_condition": "...",}\\
\texttt{~~"output\_rule": "...",}\\
\texttt{~~"exact\_values": ["literal", "tokens"],}\\
\texttt{~~"ambiguous": false}\\
\texttt{\}}

Output only \texttt{\{"constraints": [...]\}}.
\end{promptbox}

\begin{promptbox}[Fuzzy L3 Generation: Constraint Selection and Rewriting]
\textbf{\#\# Role}

You are reviewing the extracted constraint table for a coding exercise.

\textbf{\#\# Task}

Given the constraint table, do three things and return JSON:
\begin{enumerate}
\item Select at most three constraints to hide, restricted to the most unguessable literals such as deliberate typos, obscure magic constants, and non-obvious sentinel values.
\item For each hidden constraint, provide a vague replacement phrase, exact PRD strings to replace, at least five trigger keywords, and a precise oracle answer.
\item Pick one or two non-hidden constraints and rewrite them as vague internal references (\texttt{vague\_refs}).
\end{enumerate}

Do not hide action names, parameter names, default values, error categories, or feature descriptions.

Return only JSON in the form:

\texttt{\{}\\
\texttt{~~"hide": [\{...\}],}\\
\texttt{~~"vague\_refs": [\{...\}]}\\
\texttt{\}}
\end{promptbox}
